\clearpage
\fancyhead[L]{\small INSTRUMENTS.}
\SourceLeaf{415}
\begin{center}
\Large INSTRUMENTS.
\end{center}

\begin{center}
Göttingische gelehrte Anzeigen. No. 127. pp. 1257..1267. 8 August 1818.
\end{center}

\par\indent Our \emph{new observatory}, this great and eloquent monument to our government’s love of science, has so far received no special mention in our journal, although for a year and a half the building has been sufficiently complete to house the stock of instruments from the old observatory and to allow intermittent observations in one section of the building. But by their very nature these observations could only be of the kind that could already be made at the old observatory, and it did not seem appropriate to us to report on them separately in this journal. Several have already been published in astronomical journals. The new observatory, intended to be second to none, can take its rightful rank only when it has its \emph{fixed meridian instruments}, and the time when it will be fully equipped is now near. --- We shall say nothing here about the exterior of the building, which befits the dignity of its purpose and does credit to the skillful and tasteful work of our university architect, Mr. \textsc{Müller}. A full account of the interior arrangements will also be reserved for another place. But as the new principal instruments, which alone give the building its true scientific value, arrive and are set up one by one, they will be described in more detail in this journal and introduced to those who are devoted to astronomy. --- The first of the new meridian instruments, the \emph{Repsold meridian circle}, arrived in April of this year and was set up by the maker himself. The instrument, which Mr. \textsc{Repsold} of Hamburg made some time ago originally for his own use and which had stood in his private observatory on the Hamburg ramparts, is not unfamiliar to friends of astronomy, since several extracts from Mr. \textsc{Repsold}’s journals have been printed in the \textit{Monatliche Correspondenz}. We too have reported in this journal (1811, p. 1290 [p. 327 of this volume]) on some observations made with this circle by Professor \textsc{Schumacher}. Since Mr. \textsc{Repsold}’s observatory had been demolished during the period so unfortunate for Hamburg and the owner had consequently had to leave the instrument unused, our government approved Councillor \textsc{Gauss}’s proposal to purchase it for our new observatory. In doing so, the maker undertook several improvements so important and substantial that the instrument has in a sense become entirely new. The work on these new arrangements was often interrupted by unforeseen circumstances. This explains the long delay in delivering the instrument, for whose installation at the observatory everything had been ready for a year and a half. --- The design of this instrument deserves a more detailed description all the more because it is so far unique and came entirely from the ingenious maker’s own ideas. It combines a first-rate transit telescope with a circle very similar to the new Greenwich mural circle. The telescope is $7\frac{1}{2}$ Paris feet long, with a focal length of $6\frac{3}{4}$ feet and an aperture of 46 lines. It is mounted on an axis 4 feet long, which, as with other transit telescopes, is suspended between two stone piers $6\frac{1}{2}$ feet high and 22 inches square. Three separate balancing arrangements cancel out the flexure of the telescope and of the axis and bear so much of the instrument’s weight that the bearing cups are subjected to a load of only a few Loth. The bearing cups \SourceLeaf{416} themselves are made of rock crystal, and the perfectly cylindrical pivots of equal diameter are made of bell metal. If even the slight friction should wear the pivots a little after long use, a special arrangement allows them to rest on a different part of their surface. The eyepiece tube can be moved parallel to the axis to bring any of the threads into the middle of the field of view. The threads are illuminated through the axis. The highest magnification ordinarily used is 96 times. Under favorable atmospheric conditions, stars as faint as third magnitude can easily be observed in daylight. The hanging spirit level used to set the axis horizontally is made with the greatest care and allows fractions of a second to be distinguished with certainty. A single large slab serves as a common base for both piers and is itself founded on large stone blocks 12 feet deep in the ground. The stability of this arrangement has proved entirely satisfactory since the installation. The azimuth, horizontal setting of the axis, and line of collimation have shown almost no variation, at least so far. --- The circle is fixed on the axis near one end and measures $3\frac{1}{2}$ feet in diameter at the graduation. The divisions were marked with admirable accuracy by very neat lines at intervals of 5 minutes while the circle was mounted on the axis itself, so that all eccentricity was avoided. It did, of course, have to be taken off again for shipment. Refitting it on site introduced only an almost imperceptible eccentricity, which Councillor \textsc{Gauss}’s careful examination puts at $0''8$. The readings are taken through three excellent micrometer microscopes. Their supports are fixed to the movable parts of the bearings so that their distance from the face of the circle remains perfectly constant, even if that face should deviate slightly from a plane. One microscope is at the bottom, and the other two are at the sides, $90^\circ$ from it. The construction of the instrument did not permit the addition of a fourth fixed microscope (formerly there was only one). Since the two side microscopes cancel out the effect of eccentricity completely, only a third of that effect remains when all three are read. It amounts to only $0''3$ and can be allowed for wherever that seems necessary. At Councillor \textsc{Gauss}’s request, the maker will also supply a fourth, \emph{movable} microscope to examine the graduation itself. The slides in the microscopes, moved by screws whose motion is extremely uniform and whose backlash has been eliminated by \textsc{Ramsden}’s familiar invention, carry a circular opening in which the bisection of the division marks can be judged so precisely that fractions of a second can be relied on. The circle has nine spokes. Between one pair there is a copper cylinder that can be set horizontally to within half a second by another excellent spirit level. This makes any changes in the position of the microscopes apparent. It also gives absolute zenith distances, since observations with the instrument reversed give the inclination of the line of sight to this cylinder. --- When the instrument arrived, the reticle in the telescope was still the one fitted when it was first made. It consisted of three vertical spider threads and two horizontal ones $21''$ apart. When Mr. \textsc{Repsold} returned to Göttingen at the end of May, he replaced this reticle with a new one, which considerably improved the accuracy of pointing. The new reticle consists of five vertical and two horizontal spider threads selected for the purpose. The latter are perfectly parallel and are $12''7$ apart. In observing, the star is usually brought midway between these two threads. This can be done with great precision, and its passages across the vertical threads can always be observed easily and distinctly. For stars on the equator the time intervals between successive vertical threads are a little over 16 seconds. --- For familiar reasons, observations of Polaris are among the daily observations at every well-equipped \SourceLeaf{417} observatory. With a new transit telescope they are doubly important, since the principal adjustments have to be based on them. Councillor \textsc{Gauss} therefore carefully observed all the upper and lower culminations of this star whenever the very favorable June weather permitted. As an example of what the instrument can do, we give the right ascensions observed so far, together with a comparison with the auxiliary table that Mr. \textsc{Bessel} gave in the astronomical yearbook for 1817 and improved in the following year’s volume. The fraction of the day is counted from the lower culmination corresponding to the stated date. The last column gives the square of the precision of each individual result according to the number of thread crossings that contribute to it. Some astronomers have quite appropriately called this square the weight of the observation. The unit is the precision of a determination inferred from two successive culminations, each observed at only one thread. The following 19 determinations were based on 41 observed culminations, with two or three immediately successive culminations combined each time into a single result.

\begin{center}
\begin{tabular}{r@{\hspace{4pt}}l|c c c|c|c}
\multicolumn{2}{c}{1818} & \multicolumn{3}{c}{\shortstack{Observed right\\ascension}} & \shortstack{Correction to\\the table} & Weight \\
\hline
June & 3.25  & $0^{\mathrm{u}}$ & $55^{\mathrm{m}}$ & $58^{\mathrm{s}}65$ & $+3^{\mathrm{s}}93$ & 2.67 \\
     & 4.25  & & & 58.19 & $+2.80$ & 5.00 \\
     & 5.25  & & & 58.56 & $+2.45$ & 5.00 \\
     & 6.25  & & & 58.91 & $+2.22$ & 4.00 \\
     & 8.50  & $0$ & $56$ & $1.14$ & $+2.97$ & 6.40 \\
     & 10.25 & & & 2.58 & $+3.20$ & 3.75 \\
     & 11.25 & & & 2.68 & $+2.59$ & 1.67 \\
     & 12.50 & & & 3.31 & $+2.36$ & 4.29 \\
     & 16.75 & & & 6.11 & $+2.13$ & 4.44 \\
     & 18.25 & & & 7.31 & $+2.26$ & 5.00 \\
     & 19.25 & & & 8.29 & $+2.51$ & 2.86 \\
     & 21.25 & & & 10.49 & $+3.26$ & 5.00 \\
     & 22.75 & & & 12.02 & $+3.60$ & 1.33 \\
     & 26.25 & & & 13.49 & $+2\ 60$ & 1.60 \\ 
     & 27.25 & & & 14.77 & $+3.14$ & 4.44 \\
     & 29.50 & & & 16.71 & $+3.42$ & 3.43 \\
July & 7.50  & & & 22.51 & $+3.31$ & 6.00 \\
     & 8.25  & & & 22.97 & $+3.22$ & 1.67 \\
     & 9.25  & & & 23.89 & $+3.40$ & 3.75 \\
\end{tabular}
\end{center}

\par\indent The mean of the corrections to the table found from these 19 observations is $+2^{\mathrm{s}}88$. If one allows oneself to regard the table’s error as constant over this period, applying Councillor \textsc{Gauss}’s theory of probability gives a probable error $=0^{\mathrm{s}}64$ for an observation whose precision is 1, and a probable error $=0^{\mathrm{s}}076$ for the final result. As a second example we add the observations of Uranus that Councillor \textsc{Gauss} made at the time of that planet’s opposition. The declinations are based on comparison with the immediately preceding lower culmination and the immediately following upper culmination of the North Star, except on 7 and 13 June, when only the lower culmination had been observed and the declination of the North Star was taken from \textsc{Bessel}’s table. The refractions were taken from \textsc{Bessel}’s table, which is based on \textsc{Bradley}’s observations. If they were taken from \textsc{Carlini}’s table, the declinations of Uranus would come out $3''$ farther south. Once enough suitable observations have been gathered, it will become clear which of the two tables is best suited to the climate here.

\SourceLeaf{418}
\begin{center}
\setlength{\tabcolsep}{3pt}
\begin{tabular}{l|ccc|ccc|ccc}
\hline
1818 & \multicolumn{3}{c|}{Mean time} & \multicolumn{3}{c|}{Right asc.} & \multicolumn{3}{c}{South decl.} \\
\hline
June 3 & $12^{\mathrm u}$ & $22^{\mathrm m}$ & $17^{\mathrm s}5$ & $257^\circ$ & $23'$ & $28''3$ & $23^\circ$ & $1'$ & $41''5$ \\
4 & 12 & 18 & 10.9 & & 20 & 47.8 & & 1 & 29.0 \\
5 & 12 & 14 & 4.6 & & 18 & 9.9 & & 1 & 14.4 \\
7 & 12 & 5 & 51.6 & & 12 & 54.4 & & 0 & 50.5 \\
8 & 12 & 1 & 45.1 & & 10 & 15.0 & & 0 & 38.4 \\
9 & 11 & 57 & 38.6 & & 7 & 36.4 & & 0 & 27.4 \\
10 & 11 & 53 & 32.1 & & 4 & 56.5 & & 0 & 12.6 \\
11 & 11 & 49 & 25.7 & & 2 & 18.0 & & 0 & 2.9 \\
12 & 11 & 45 & 19.2 & 256 & 59 & 33.9 & 22 & 59 & $50\;9$ \\ 
13 & 11 & 41 & 12.7 & & 57 & 0.4 & & 59 & 40.5 \\
\hline
\end{tabular}
\end{center}

\Needspace{18\baselineskip}\par\indent Councillor \textsc{Gauss} entrusted the comparison of these observations with \textsc{Delambre}’s tables of Uranus to Mr. \textsc{Dirksen}, who pursues the study of astronomy here with exceptional zeal. The results of this comparison are as follows:

\begin{center}
\begin{tabular}{c|cc}
& \multicolumn{2}{c}{Difference} \\
& Right asc. & Declination \\
\hline
June 3 & $-54.1$ & $+10.3$ \\
4 & $-52.1$ & $+10.6$ \\
5 & $-52.9$ & $+12.9$ \\
7 & $-55.8$ & $+12.3$ \\
8 & $-55.9$ & $+12.0$ \\
9 & $-56\;8$ & $+10.5$ \\
10 & $-56.4$ & $+12.9$ \\
11 & $-57.2$ & $+10.1$ \\
12 & $-52.4$ & $+9\;4$ \\ 
13 & $-58.1$ & $+7.4$ \\
\hline
Mean & $-55.2$ & $+10.8$ \\
\end{tabular}
\end{center}

\Needspace{15\baselineskip}\par\indent From this, Mr. \textsc{Dirksen}’s calculation gave:

\begin{center}
Time of opposition, 9 June 1818.\\
$5^{\mathrm u}\,30^{\mathrm m}\,43^{\mathrm s}$ mean time in Göttingen
\end{center}
\begin{center}
\begin{tabular}{ll}
True longitude . . . . . . . . . . . . . . & $258^\circ\;10'\;34''5$ \\
Geocentric latitude . . . . . . . . . . . . & $0\;4\;12.1$ \\
Heliocentric . . . . . . . . . . . . . . & $0\;3\;58.7$ South \\
\end{tabular}
\end{center}

\noindent and the error in \textsc{Delambre}’s tables

\begin{center}
\begin{tabular}{ll}
in longitude . . . . . . . . . . . . . . & $-45''3$ \\
in heliocentric latitude . . . . . . . . . . & $+14.6$ \\
\end{tabular}
\end{center}

\par\indent We cannot close this report without also mentioning \textsc{Shelton}’s pendulum clock. It had already been in use for 47 years at the old observatory, but its excellence can only now be fully appreciated, since the observatory previously had no means of determining the time with the utmost precision. By way of evidence we give only the record of its rate over one month, and note that the same uniformity held throughout the time when the \textsc{Repsold} instrument was used to determine its rate. The date refers to solar days. To make the table easier to read, the hours shown by the clock continue past 24 in each case, up to the following noon. The last column gives the number of stars from \textsc{Maskelyne}’s fundamental catalogue on which the time determination is based.

\SourceLeaf{419}
\begin{center}
\begin{tabular}{ll|c|c|c}
 & & Clock time & \shortstack{Offset from\\sidereal time} & \shortstack{Stars\\compared} \\
\hline
June & 15 & $15^{\mathrm u}\;22^{\mathrm m}$ & $+0^{\mathrm s}41$ & 4 \\
 & 17 & $20\;\;50$ & $+0.28$ & 2 \\
 & 18 & $11\;\;37$ & $+0.29$ & 2 \\
 & 20 & $12\;\;28$ & $+0.11$ & 2 \\
 & 21 & $11\;\;37$ & $-0.06$ & 2 \\
 & 22 & $7\;\;\;30$ & $-0.59$ & 1 \\
 & 23 & $11\;\;40$ & $-0.90$ & 1 \\
 & 26 & $14\;\;54$ & $-1.78$ & 5 \\
 & 27 & $16\;\;15$ & $-1.86$ & 4 \\
 & 28 & $14\;\;41$ & $-1.99$ & 1 \\
 & 29 & $18\;\;\;\;1$ & $-2.32$ & 6 \\
 & 30 & $16\;\;\;\;9$ & $-2.43$ & 5 \\
\end{tabular}
\end{center}

\begin{center}
\begin{tabular}{ll|c|c|c}
 & & Clock time & \shortstack{Offset from\\sidereal time} & \shortstack{Stars\\compared} \\
\hline
July & 1 & $14^{\mathrm u}\;7^{\mathrm m}$ & $-2^{\mathrm s}92$ & 1 \\
 & 2 & $14\;\;51$ & $-3.11$ & 2 \\
 & 6 & $25\;\;17$ & $-3.87$ & 4 \\
 & 7 & $22\;\;31$ & $-3.90$ & 7 \\
 & 8 & $16\;\;36$ & $-4.03$ & 4 \\
 & 9 & $28\;\;\;7$ & $-4.21$ & 4 \\
 & 10 & $15\;\;\;3$ & $-4.28$ & 3 \\
 & 13 & $29\;\;24$ & $-4.53$ & 2 \\
 & 14 & $29\;\;14$ & $-4.62$ & 4 \\
 & 15 & $15\;\;56$ & $-4.73$ & 5 \\
 & 16 & $14\;\;41$ & $-4.93$ & 1 \\
\end{tabular}
\end{center}

\par\indent It should also be mentioned that the clock’s installation is still provisional. It rests on a separate stone foundation, supported only by its case. It has no direct attachment to the stone pier beside the case. It runs for a month on a single winding. About halfway through that time, the weight comes opposite the pendulum bob. This had sometimes been considered detrimental, and attempts had been made to avoid it by various arrangements. Yet the record for \textsc{Shelton}’s clock shows no perceptible change in rate at that time. Indeed, its rate is much more regular than that of many other clocks at the foremost observatories, including clocks attached directly to a stone pier. Should we not conclude from this that both circumstances are of much less importance than has generally been thought up to now?
